\documentclass[11pt]{article}
\usepackage[margin=1in]{geometry}
\usepackage{amsmath,amssymb,amsthm,mathtools}
\usepackage{graphicx}
\usepackage{microtype}
\usepackage[hidelinks]{hyperref}

\newtheorem{theorem}{Theorem}
\newcommand{\dom}{\operatorname{dom}}
\newcommand{\Ktilde}{\widetilde{\mathcal K}}
\newcommand{\I}{\mathcal I}

\title{A Continuum of Evolution-Generator Extensions\\
of One Evolution Pre-Generator}
\author{Automated mathematics run \texttt{fa\_banach\_001}, lane 1}
\date{August 9, 2026}

\begin{document}
\maketitle

\begin{center}
\fbox{\parbox{0.9\textwidth}{\textbf{Status.}
Candidate full solution, likely valid, pending expert review.  A single
canonical evolution pre-generator associated with a constant closed operator
on $L^2(0,1)$ has continuum many pairwise distinct evolution-generator
extensions.  Equivalently, the corresponding correctly posed problem has
continuum many solution operators in the weak sense of the source paper.}}
\end{center}

\section{Source question}

Neidhardt, Stephan, and Zagrebnov~\cite{NSZ} work on
$L^p(\I,X)$, where $\I=[0,T]$, $T>0$, $1\leq p<\infty$, and $X$ is a
separable Banach space.  Their Definition~3.1 calls an operator $\mathcal K$
an \emph{evolution pre-generator} if its domain consists of continuous
functions, is invariant under multiplication by $W^{1,\infty}(\I)$ scalar
functions, satisfies
\[
 \mathcal K M(\phi)f-M(\phi)\mathcal Kf=M(\dot\phi)f,
\]
and has a dense cross-section in $X$ at every $t\in(0,T]$.  It is an
\emph{evolution generator} if, in addition, it generates a strongly
continuous semigroup under the paper's sign convention $e^{-\tau\mathcal K}$.

For a constant closed operator $C$ on $X$, let $\mathcal C$ denote its
induced multiplication operator on $L^p(\I,X)$, and let
\[
 D_0 f=f',\qquad
 \dom(D_0)=\{f\in W^{1,p}(\I,X):f(0)=0\}.
\]
The canonical evolution pre-generator is
\begin{equation}\label{eq:pregenerator}
 \Ktilde=D_0+\mathcal C,
 \qquad \dom(\Ktilde)=\dom(D_0)\cap\dom(\mathcal C).
\end{equation}
In Remark~3.6(i), on PDF page 14 (printed page 12), the paper asks:
\begin{quote}
``It is an open problem whether an evolution pre-generator admits several
extension which are evolition generators.''
\end{quote}
The spelling is transcribed exactly.  The remark adds that such an example
would give a non-autonomous Cauchy problem with more than one solution
operator under Definition~3.5.

\begin{figure}[p]
  \centering
  \includegraphics[width=\textwidth]{figures/open_problem_crop.png}
  \caption{Definition~3.5, Remark~3.6, and the immediate uniqueness context
  in Neidhardt--Stephan--Zagrebnov~\cite{NSZ}, PDF page 14.}
\end{figure}
\clearpage

\section{Proof intuition}

The question is an extension problem, so begin with a familiar closed
operator having many semigroup-generator extensions.  The closed minimal
Laplacian on $(0,1)$ is contained in every Robin Laplacian.  Each nonnegative
Robin realization is self-adjoint and hence generates a contraction heat
semigroup.  Passing from a spatial heat semigroup to the associated evolution
semigroup on $L^p([0,T],L^2(0,1))$ turns every Robin extension into an
evolution-generator extension of the same canonical pre-generator.

The dense-cross-section requirement does not destroy this construction.  At
any positive time, a smooth scalar cutoff which vanishes at time zero can be
multiplied by an arbitrary vector from the dense domain of the minimal
Laplacian.  Finally, different Robin parameters impose different boundary
conditions, so their spatial generators, propagators, and evolution
generators are pairwise distinct.  Since the Robin parameter ranges over an
interval, the one pre-generator has continuum many extensions.

\section{Full answer}

Let $H=L^2(0,1)$.  Define the closed minimal Laplacian
\begin{equation}\label{eq:minimal}
 Su=-u'',\qquad
 \dom(S)=H^2_0(0,1)
 =\{u\in H^2(0,1):u(0)=u'(0)=u(1)=u'(1)=0\}.
\end{equation}
Equivalently, $S$ is the closure of $-d^2/dx^2$ initially defined on
$C_c^\infty(0,1)$.

\begin{theorem}\label{thm:main}
Fix $T>0$ and $1\leq p<\infty$.  Let $\mathcal S$ be the multiplication
operator on $L^p([0,T],H)$ induced by the closed minimal Laplacian $S$, and set
\[
 \Ktilde=D_0+\mathcal S,
 \qquad \dom(\Ktilde)=\dom(D_0)\cap\dom(\mathcal S).
\]
Then $\Ktilde$ is an evolution pre-generator.  Moreover, it has continuum
many pairwise distinct closed extensions which are evolution generators.
Consequently, the constant problem
\[
 \dot u(t)=-Su(t)
\]
has continuum many solution operators in the sense of
Definition~3.5(ii) of~\cite{NSZ}.
\end{theorem}

\begin{proof}
The operator $S$ is closed and densely defined on the separable Hilbert space
$H$.  Its induced multiplication operator is
\[
 \dom(\mathcal S)=
 \{f\in L^p(\I,H): f(t)\in\dom(S)\text{ a.e. and }Sf(\cdot)\in L^p(\I,H)\}.
\]
We first verify the pre-generator property.  Since
$\dom(\Ktilde)\subset\dom(D_0)$, every domain element has a continuous
representative.  For $\phi\in W^{1,\infty}(\I)$, scalar multiplication
preserves both $\dom(D_0)$ and $\dom(\mathcal S)$, and the Leibniz rule gives
\[
 \Ktilde M(\phi)f-M(\phi)\Ktilde f=M(\dot\phi)f.
\]
It remains to check the dense cross-section.  Fix $t_0\in(0,T]$, $x\in H$,
and $\varepsilon>0$.  Because $\dom(S)$ is dense in $H$, choose
$u\in\dom(S)$ with $\|u-x\|_H<\varepsilon$.  Choose
$\psi\in C^\infty([0,T])$ with $\psi(0)=0$ and $\psi(t_0)=1$.  Then
$f(t)=\psi(t)u$ lies in $\dom(\Ktilde)$ and
$\|f(t_0)-x\|_H<\varepsilon$.  Thus every positive-time cross-section is
dense.  This also recovers directly the hypotheses of Lemma~3.4
of~\cite{NSZ}.

For each $a\in[0,\infty)$, consider the densely defined nonnegative closed
form
\begin{equation}\label{eq:form}
 q_a(u,v)=\int_0^1 u'(x)\overline{v'(x)}\,dx
 +a\bigl(u(0)\overline{v(0)}+u(1)\overline{v(1)}\bigr),
 \qquad \dom(q_a)=H^1(0,1).
\end{equation}
The first representation theorem for closed forms~\cite[Chapter~VI]{Kato}
gives a nonnegative self-adjoint operator $A_a$.  Integration by parts
identifies it as the Robin Laplacian
\begin{equation}\label{eq:robin}
 A_a u=-u'',\qquad
 \dom(A_a)=\{u\in H^2(0,1):u'(0)=a u(0),\ u'(1)=-a u(1)\}.
\end{equation}
Every $u\in\dom(S)$ satisfies these boundary conditions, and the two
operators agree there.  Hence $S\subset A_a$ for every $a\geq0$.

Since $A_a$ is nonnegative self-adjoint, $\{e^{-rA_a}\}_{r\geq0}$ is a
strongly continuous contraction semigroup on $H$.  It defines the propagator
\[
 U_a(t,s)=e^{-(t-s)A_a},\qquad 0<s\leq t\leq T,
\]
and the evolution semigroup on $L^p(\I,H)$
\begin{equation}\label{eq:evolution}
 (e^{-\tau\mathcal K_a}f)(t)
 =e^{-\tau A_a}f(t-\tau)\,\mathbf 1_{\I}(t-\tau).
\end{equation}
Let $\mathcal K_a$ denote its generator in the sign convention of the
source.  Theorem~3.3 of~\cite{NSZ} makes $\mathcal K_a$ an evolution
generator.  Alternatively, equation~\eqref{eq:evolution} is exactly the
defining evolution-semigroup formula.

Let $\mathcal A_a$ be the multiplication operator induced by $A_a$.
Lemma~3.8(ii) of~\cite{NSZ} states that
\[
 \mathcal K_a f=D_0f+\mathcal A_a f,
 \qquad f\in\dom(D_0)\cap\dom(\mathcal A_a).
\]
The inclusion $S\subset A_a$ implies
$\dom(\mathcal S)\subset\dom(\mathcal A_a)$ and
$\mathcal A_a f=\mathcal S f$ on $\dom(\mathcal S)$.  Therefore
\[
 \Ktilde\subset\mathcal K_a
 \qquad(a\geq0).
\]
Thus every $\mathcal K_a$ is a closed evolution-generator extension of the
same evolution pre-generator $\Ktilde$.

It remains only to prove that these extensions are distinct.  If
$a\ne b$, take the cubic Hermite interpolant $u\in H^2(0,1)$ determined by
\[
 u(0)=1,\quad u'(0)=a,\quad u(1)=0,\quad u'(1)=0.
\]
Then $u\in\dom(A_a)$ but $u\notin\dom(A_b)$, so $A_a\ne A_b$.  If
$\mathcal K_a=\mathcal K_b$, their evolution semigroups would coincide.
By the one-to-one correspondence of Theorem~3.3 of~\cite{NSZ}, their
propagators would coincide, so
$e^{-rA_a}=e^{-rA_b}$ for every $0<r\leq T$.  Taking the strong derivative at
$r=0$ would give $A_a=A_b$, a contradiction.  Hence the $\mathcal K_a$ are
pairwise distinct.  Since $[0,\infty)$ has cardinality continuum, the proof is
complete.
\end{proof}

\section{Verification, scope, and novelty}

The construction stays inside the exact framework of the source: $H$ is
separable; $S$ is closed and densely defined; $\Ktilde$ has the canonical form
\eqref{eq:pregenerator}; and every extension comes from a bounded strongly
continuous propagator.  The source itself emphasizes that its notion of
solution operator is weak.  The theorem answers that precise weak question;
it does not claim classical well-posedness for the minimal-Laplacian equation.
Indeed, Robin heat trajectories need not remain in $\dom(S)$.

No computation is used in the proof.  The two key checks are (1) that
$S\subset A_a$ makes the full canonical domain of $D_0+\mathcal S$, not just
an algebraic tensor core, lie in every $\mathcal K_a$, and (2) that different
Robin boundary conditions force different propagators and hence different
evolution generators.  The detailed audit is in \path{VERIFICATION.md}.

The local run indexes were searched for arXiv:1612.06147, ``evolution
pre-generator,'' multiple or several evolution-generator extensions,
nonunique solution operators, the minimal Laplacian, and Robin extensions.
Bounded web searches on August 9, 2026 used the exact open-problem sentence,
the arXiv id and DOI, and combinations of those terms.  Exact-phrase searches
returned the source paper; searches combining the terminology with Laplacian,
Robin, and nonuniqueness found no later resolution.  A citation-metadata check
also inspected the titles of the small related Trotter/evolution-semigroup
citation set and found no stated resolution.  This is a bounded search, not a
systematic bibliographic proof of novelty, so expert literature review remains
necessary.  Full search bounds are recorded in \path{novelty_search.md}.

\begin{thebibliography}{9}
\bibitem{NSZ}
H. Neidhardt, A. Stephan, and V. A. Zagrebnov,
\emph{Convergence rate estimates for Trotter product approximations of
solution operators for non-autonomous Cauchy problems},
Publ. Res. Inst. Math. Sci. \textbf{56} (2020), no.~1, 83--135;
arXiv:1612.06147; doi:10.4171/PRIMS/56-1-5.

\bibitem{Kato}
T. Kato,
\emph{Perturbation Theory for Linear Operators}, second edition,
Springer, Berlin, 1976.
\end{thebibliography}

\end{document}
